% 18b-lineage.tex: what Polaris inherits, what it combines, and what may be its own.
% Added 2026-09-25, after an outside reader observed that the report's internal traceability was
% meticulous and its academic ancestry faint: until this section it cited nothing.
\section{Lineage: what is inherited, and what is not}
\label{sec:lineage}

\motif{Putting known pieces together is not inventing them.}

Until this section the report cited nothing, and a reader could not tell invention from synthesis.
This section separates them. It claims less than a reader might assume: Polaris contains no new
cryptography, no new protocol and no new proof system. Every primitive below is someone else's, and
what follows names whose.

\subsection{Inherited}

\begin{itemize}
  \item \emph{Signatures.} ML-DSA and SLH-DSA are the NIST standards \cite{fips204,fips205}. The
    migration path and its audit trail are Polaris's engineering; the security of the primitives is
    not, and nothing in this repository bears on it.
  \item \emph{Merkle commitments and transparency.} The epoch tree and the anchor batches are Merkle
    trees \cite{merkle1988}. The transparency logs follow Certificate Transparency's append-only log,
    inclusion and consistency proofs \cite{rfc6962}, and the witness cosignature over a log head
    follows the idea of collective witness cosigning \cite{syta2016}.
  \item \emph{Anonymous membership.} Proving membership of a committed set without revealing which
    member, with a nullifier to refuse a second use, is the construction Zerocash made practical
    \cite{zerocash2014}; the circuit is built on Plonky2 \cite{plonky2}. The broader aim, credentials
    a holder shows without being traced, goes back to Chaum \cite{chaum1985} and to anonymous
    credential systems \cite{camlys2001}, which Polaris does not implement: it is not a general
    anonymous-credential system (Section~\ref{sec:privacy}).
  \item \emph{Presentation and wallets.} The wallet-facing verifier implements OpenID for Verifiable
    Presentations and its high-assurance profile \cite{oid4vp,haip}, over SD-JWT and SD-JWT VC
    \cite{sdjwt,sdjwtvc} with the token status list \cite{statuslist}. The mobile-document route
    follows ISO/IEC 18013-5 \cite{iso18013}, and the credential vocabulary the W3C data model
    \cite{vcdm2}.
  \item \emph{Relying parties.} Per-relying-party subject identifiers are OpenID Connect's pairwise
    identifiers \cite{oidc}. The short-lived signed status assertion is the pattern of an OCSP
    response a server staples \cite{rfc6960}. Operator sign-in uses WebAuthn \cite{webauthn2}.
  \item \emph{Proofing.} The evidence strengths and the derivation of an assurance level follow NIST
    SP 800-63A \cite{sp80063a}; Polaris records the reading, and says so where it is a reading.
  \item \emph{Duress codes.} A second secret that signals coercion is long-standing practice in
    alarm systems and banking; it is not invented here. What Polaris adds is the measurement in
    Section~\ref{sec:limitations} of the adversary against whom it is worse than nothing.
  \item \emph{Testing the tests.} Mutation analysis is DeMillo, Lipton and Sayward's
    \cite{demillo1978}, surveyed by Jia and Harman \cite{jia2011}.
\end{itemize}

\subsection{What may be Polaris's own}

To the author's knowledge, and offered as a description, not a claim of priority:
\begin{itemize}
  \item the \emph{arrangement}: rights-level constraints held by the schema rather than the
    application, each paired with a machine check that has its own detection test;
  \item the \emph{assurance method}: mutating the assurance layer itself (the checks, triggers,
    procedures, verifiers and witnesses) to find green checks that prove the wrong proposition, as
    described in Section~\ref{sec:senses}, with its record in \evidencerecord;
  \item the \emph{operating contract}: progress counted only as a surviving external dependency,
    with a kill line, rather than as internal growth.
\end{itemize}
None of these is a result a reader should accept on this report's word. Each is described so that it
can be compared with prior work this section does not yet cite, and a reader who knows such work is
asked to name it.

\subsection{What this section does not settle}

It does not survey identity systems, digital-identity programmes or the anonymous-credential
literature; it names only what Polaris's mechanisms directly derive from. A comparison with deployed
national systems is the table in the repository's README, which states that every one of them is
deployed and Polaris is not.
